\documentclass[11pt, a4paper]{article}
\usepackage[a4paper, margin=2.4cm]{geometry}
\usepackage{fontspec}
\usepackage[chinese, bidi=basic, provide=*]{babel}
\babelprovide[import, onchar=ids fonts]{english}
\babelfont{rm}{Noto Sans}
\babelfont[chinese]{rm}{Noto Sans CJK SC}
\usepackage{enumitem}
\setlist[itemize]{label=-}
\usepackage{amsmath, amssymb, booktabs, xcolor}
\usepackage[hidelinks]{hyperref}

\definecolor{accent}{RGB}{24, 72, 130}
\definecolor{muted}{RGB}{90, 98, 110}

\title{\textbf{\color{accent}基于线性规划松弛与区间对偶认证的奇完全数终局下界与路障剪枝研究}\\[0.4em]
\large 路线 A 阶段性技术报告}
\author{研究整理报告}
\date{}

\begin{document}
\maketitle

\begin{abstract}
在 Ochem 与 Rao（2012）建立的奇完全数（Odd Perfect Number, OPN）下界证明框架中，无条件规模下界 $N > B$ 由前半段“因子链与路障绕过”以及后半段“无禁用素数因子的终局（End game）下界”两部分构成。尽管最新未发表记录已通过持续扩充大整数分解库与禁用素数集将无条件下界推至 $10^{2200}$，但原框架在终局估计与路障绕过上仍沿用较为粗放的贪心乘积界。本报告总结了路线 A 的核心算法改进：通过引入欧拉唯一特殊素数约束、精确有限次幂丰度函数 $\sigma_{-1}(p^a)$ 以及 $\sigma(p^a)$ 的二次素因子蕴含约束，构建了覆盖无穷素数尾部的严格线性规划（LP）松弛模型，并结合高精度区间算术对偶认证消除浮点误差。实验表明：（1）在完全沿用 2012 年论文 15 个禁用素数集 $S_{2012}$ 的条件下，严格认证的终局下界由 $10^{1735}$ 跃升至 $10^{3129.458561}$（浮点间隙仅 $7.77 \times 10^{-8}$）；（2）在仅使用 8 个经典素数 $S_{\mathrm{BCR}}$ 时，终局下界即可达到 $10^{1725.52}$；（3）将该模型移植至论文图 1 的嵌套路障案例 Roadblock 2，可将需要分支的素数上界由 $2625$（约 $376$ 个素数）压低至 $859$（仅 $143$ 个素数），单层路障分支宽度缩减 $62\%$。这些结果为后续路线 B 基于现有分解库冲击 $10^{2200}$ 以上无条件下界提供了显著的算法冗余。
\end{abstract}

\section{研究背景与问题定义}

设 $N$ 为奇完全数，记 $\sigma(n)$ 为正整数 $n$ 的所有正因子之和，$\sigma_{-1}(n) = \sigma(n)/n$ 为 $n$ 的丰度（abundancy），则 $N$ 为完全数当且仅当 $\sigma_{-1}(N) = 2$。由 Euler 定理，任何奇完全数均可唯一表为
\[
N = p^e m^2, \qquad p \equiv e \equiv 1 \pmod{4}, \quad \gcd(p, m) = 1,
\]
其中 $p$ 称为欧拉特殊素数（special prime），其余素因子的重数均为偶数。

Ochem 与 Rao 在 2012 年发表的论文中证明了正式发表的最高无条件下界 $N > 10^{1500}$，近年 Ochem 沿用同一证明树框架与持续扩充的协作分解库（\texttt{checkfacts.txt.gz}）将无条件下界推进至 $10^{2200}$（下一步目标为 $10^{2300}$）。该框架的证明由前后两部分拼接而成：
\begin{itemize}
    \item \textbf{前半段（因子链与路障绕过）}：选定有限禁用素数集 $S$，通过递归分支假设 $p^a \parallel N$ 并利用 $\sigma(p^a) \mid 2N$ 推导矛盾，证明若 $N$ 含有 $S$ 中任意素因子则 $N > B$。当已知素因子均已分支且剩余均为未分解合数时，形成路障（Roadblock），需计算新素数的枚举上界 $b$ 并逐一分支绕过。
    \item \textbf{后半段（终局下界 / End game）}：证明若 $\gcd\bigl(N, \prod_{s \in S} s\bigr) = 1$，则为了使总丰度达到 $2$，必然有 $N > B$。
\end{itemize}

审视 Ochem 与 Rao（2012）对 15 个禁用素数
\[
S_{2012} = \{127, 19, 7, 11, 331, 31, 97, 61, 13, 398581, 1093, 3, 5, 307, 17\}
\]
所做的终局论证（给出 $10^{1735}$）及路障绕过公式，存在三处显著的数学松弛：
\begin{enumerate}
    \item \textbf{未施加欧拉特殊素数的唯一性约束}：原方法对所有允许的素数仅要求 $4 \nmid \sigma(p^a)$ 且 $\sigma(p^a)$ 不含 $S$ 中因子，导致所有 $p \equiv 1 \pmod{4}$ 且 $(p+1)/2$ 不含 $S$ 中因子的素数均被允许取 $a=1$。事实上，至多只有一个素数能取 $a \equiv 1 \pmod{4}$，其余所有 $p \equiv 1 \pmod{12}$ 的素数因 $3 \mid \sigma(p^2)$ 必须跳至 $a \ge 4$。
    \item \textbf{规模取有限幂而丰度取无穷幂}：原方法在计算规模下界时使用最小合法幂次 $p_i^{a_i}$，但在计算丰度上界时直接取 $\sigma_{-1}(p_i^{\infty}) = p_i/(p_i-1)$，造成丰度高估。
    \item \textbf{忽略 $\sigma(p^a)$ 的二次素因子强迫效应}：若 $p^a \parallel N$ 且奇素数幂 $q^k \parallel \sigma(p^a)$，则 $q^k \mid N$ 必须成立。例如对大量小素数平方 $p^2$，其 $\sigma(p^2) = p^2+p+1$ 常为素数或含大素因子 $q$，迫使 $N$ 同时包含 $q^2 \approx p^4$（除非 $q$ 恰好为唯一特殊素数），从而使 $p^2$ 的真实对数代价成倍上升。
\end{enumerate}

\section{严格线性规划松弛模型与区间对偶认证}

\subsection{覆盖无穷素数尾部的通用 LP 松弛}

为确保所得下界无任何截断漏洞，我们设计了一个同时适用于“终局下界”与“路障绕过”的通用线性规划模型。设 $S_{\mathrm{forb}}$ 为禁止出现在 $N$ 中的素数集，$E$ 为当前已知整除 $N$ 的素数集（终局情形下 $E = \varnothing$），$p_{\min}$ 为允许引入的新素数下界（终局情形下 $p_{\min} = 3$）。注意到所有满足 $q < p_{\min}$ 且 $q \notin E$ 的奇素数亦不得整除 $N$，故自动并入有效禁用集 $S_{\mathrm{eff}}$。

选定显式截断素数 $P_{\max}$ 与指数截断阈值 $a_{\mathrm{cut}}$，将所有可用奇素数划分为：
\begin{itemize}
    \item \textbf{显式素数池 $Q$}：包含所有满足 $p_{\min} \le p \le P_{\max}$ 且 $p \notin S_{\mathrm{eff}} \cup E$ 的素数，以及在分解基础池素数 $\sigma(p^a)$（满足 $p^a < F_{\mathrm{lim}}$）时产生的所有外溢奇素因子 $q > P_{\max}$（且 $q \notin E$）。
    \item \textbf{无穷尾素数}：所有不属于 $Q \cup S_{\mathrm{eff}} \cup E$ 的素数（均严格大于 $P_{\max}$）。
\end{itemize}

对每个 $p \in Q$，枚举其所有小于 $a_{\mathrm{cut}}$ 的合法偶数幂次 $a \in A_{\mathrm{even}}(p)$（满足 $S_{\mathrm{eff}}$ 中无素数整除 $\sigma(p^a)$），并添加一个偶数高次幂打包项（代表所有 $a \ge a_{\mathrm{cut}}$，对数代价保守取 $a_{\mathrm{cut}}\ln p$，丰度保守取 $\ln\frac{p}{p-1}$）。若允许引入特殊素数且 $p \equiv 1 \pmod{4}$，同理枚举合法特殊幂次 $a \in A_{\mathrm{spec}}(p)$（$a \equiv 1 \pmod{4}$）及其高次幂打包项。记 $A(p) = A_{\mathrm{even}}(p) \cup A_{\mathrm{spec}}(p)$，引入松弛变量 $x_{p,a} \in [0, 1]$。

对于无穷尾部的非特殊素数，其任意合法偶数幂次 $a \ge 2$ 的丰度效率严格小于
\[
\eta_{\mathrm{tail}} = \frac{\ln\bigl(P_{\max}/(P_{\max}-1)\bigr)}{2\ln P_{\max}}.
\]
因此引入非负聚合变量 $y_{\mathrm{tail}} \ge 0$ 表示所有尾部非特殊素数贡献的总对数丰度，其最小对数规模代价为 $y_{\mathrm{tail}} / \eta_{\mathrm{tail}}$；若允许特殊素数落入尾部，再引入变量 $x_{\mathrm{tail\_spec}} \in [0, 1]$，代价下界为 $\ln P_{\max}$，丰度上界为 $\ln\frac{P_{\max}}{P_{\max}-1}$。

记目标丰度比值为 $R$（终局时 $R = 2$；路障时 $R = 2/a'$），最小化 $\ln N$（或剩余部分 $\ln(N/M)$）的线性规划为：
\begin{equation}
\min \sum_{p \in Q}\sum_{a \in A(p)} (a\ln p)\, x_{p,a} + \frac{1}{\eta_{\mathrm{tail}}}\, y_{\mathrm{tail}} + (\ln P_{\max})\, x_{\mathrm{tail\_spec}}
\end{equation}
满足如下四组线性约束：
\begin{align}
\sum_{p \in Q}\sum_{a \in A(p)} \ln\sigma_{-1}(p^a)\, x_{p,a} + y_{\mathrm{tail}} + \ln\frac{P_{\max}}{P_{\max}-1}\, x_{\mathrm{tail\_spec}} &\ge \ln R, \\
\sum_{p \in Q,\, p \equiv 1 \pmod{4}} \sum_{a \in A_{\mathrm{spec}}(p)} x_{p,a} + x_{\mathrm{tail\_spec}} &\le 1, \\
\sum_{a \in A(p)} x_{p,a} &\le 1, \qquad \forall p \in Q, \\
x_{p,a} - \sum_{\substack{b \in A'(q) \\ b \ge k}} x_{q,b} &\le 0, \qquad \forall (p,a,q,k) \in \mathcal{I}.
\end{align}
其中 $\mathcal{I}$ 为分解 $\sigma(p^a)$ 得到的已知二次奇素因子幂 $q^k \parallel \sigma(p^a)$（$q \notin E$）集合；当 $a \equiv 1 \pmod{4}$ 时，由于 $p$ 已占用唯一特殊素数，式 (5) 中的合法幂次集 $A'(q)$ 进一步收紧为 $A_{\mathrm{even}}(q)$（并包含偶数高次幂打包项）。

\subsection{高精度区间算术对偶认证}

数值 LP 求解器（如 HiGHS）工作在双精度浮点下，其原始目标值不能直接作为严格数学定理。将上述 LP 记为标准不等式形式 $\min \{c^T x \mid A x \le b,\; 0 \le x_i \le u_i\}$（其中仅 $y_{\mathrm{tail}}$ 的上界 $u_{\mathrm{tail}} = +\infty$，其余变量 $u_i = 1$）。由线性规划弱对偶定理，对任意非负对偶向量 $\lambda \ge 0$，只要满足 $\lambda_0 \le c_{\mathrm{tail}} = 1/\eta_{\mathrm{tail}}$（保证无界变量 $y_{\mathrm{tail}}$ 的检验数 $c_{\mathrm{tail}} - \lambda_0 \ge 0$），则对任意原始可行解 $x$ 恒有：
\begin{equation}
c^T x \ge -\lambda^T b + \sum_{u_i = 1} \min\Bigl(0,\; c_i + (A^T \lambda)_i\Bigr).
\end{equation}
具体认证流程为：先由 HiGHS 求解获得对偶乘子近似值 $\lambda \ge 0$；在 30 位十进制精度的区间算术（\texttt{mpmath.iv}, \texttt{dps=30}）下，将 $\lambda_0$ 严格截断至 $c_{\mathrm{tail}}$ 的区间下端点以保证无界检验数非负，随后全部采用区间算术计算式 (6) 右端并取下端点除以 $\ln 10$ 的区间上端点。所得数值为 100\% 无浮点漏洞的严格数学下界。

\section{终局下界实验与禁用素数集消融分析}

实验参数统一设置为 $P_{\max} = 4000$、$a_{\mathrm{cut}} = 16$、二次分解上限 $F_{\mathrm{lim}} = 10^{14}$，使用 Python 结合 \texttt{sympy}、\texttt{scipy.optimize.linprog}（HiGHS）与 \texttt{mpmath.iv} 实现。

\subsection{主实验：$S_{2012}$ 上的分步增益与严格认证}

表 \ref{tab:main_endgame} 给出了在完全相同的 15 个禁用素数集 $S_{2012}$ 下，各模型给出的终局规模下界对比。

\begin{table}[htbp]
\centering
\caption{$S_{2012}$（15 个禁用素数）上的终局规模下界对比}
\label{tab:main_endgame}
\begin{tabular}{@{}lcc@{}}
\toprule
方法 / 模型 & 终局下界（$\log_{10} N$） & 相对基线增益 \\
\midrule
Ochem \& Rao (2012) 贪心乘积法 & $1735$ & 基线 \\
Model A（唯一特殊素数 + 精确 $\sigma_{-1}(p^a)$） & $1995.75$ & $+260.75$ \\
Model B（Model A + 二次因子蕴含约束，浮点 LP） & $3129.458561$ & $+1394.46$ \\
\textbf{Model B 区间算术对偶认证（\texttt{mpmath.iv}, 30 dps）} & $\mathbf{3129.458561}$ & $\mathbf{+1394.46}$ \\
\bottomrule
\end{tabular}
\end{table}

区间对偶下界与浮点 LP 解的差值仅为 $7.77 \times 10^{-8}$，验证了弱对偶证书的高度紧致性。由此可确立如下严格命题：

\medskip
\noindent\textbf{命题 1（严格认证结果）} \textit{若奇完全数 $N$ 不被 $S_{2012} = \{3, 5, 7, 11, 13, 17, 19, 31, 61, 97, 127, 307, 331, 1093, 398581\}$ 中的任何素数整除，则 $N > 10^{3129.458561}$。}
\medskip

从表 \ref{tab:main_endgame} 可见，Model B 相较 Model A 暴涨超过 $1133$ 个数量级。核心原因在于：在无二次蕴含约束时，LP 倾向于大量选取效率最高的小素数平方 $p^2$（对数代价仅 $2\ln p$）；引入式 (5) 后，由于 $\sigma(p^2) = p^2 + p + 1$ 往往包含与 $p^2$ 同量级的大素因子 $q$，且至多只有一个 $q$ 可作特殊素数，其余均须至少以 $q^2 \approx p^4$ 进入 $N$，从而使 $p^2$ 的有效代价从 $2\ln p$ 激增至约 $6\ln p$，迫使最优解转向更高幂次或更大素数。

\subsection{禁用素数集留一法（Leave-One-Out）消融实验}

为评估 $S_{2012}$ 中每个素数对终局下界的边际贡献，并测试反向压缩禁用集的可行性，我们对 $S_{2012}$ 逐一剔除单个素数，并同步测试经典 8 素数集 $S_{\mathrm{BCR}} = \{3, 5, 7, 11, 13, 19, 31, 127\}$。32 项并行 LP 任务（含后文 Roadblock 2 扫描）总耗时仅 $6.16$ 秒，消融结果汇总于表 \ref{tab:ablation}。

\begin{table}[htbp]
\centering
\caption{禁用素数集消融实验结果（Model B，$P_{\max}=4000, a_{\mathrm{cut}}=16$）}
\label{tab:ablation}
\begin{tabular}{@{}lcc|lcc@{}}
\toprule
禁用素数集 & $|S|$ & $\log_{10} N$ 下界 & 禁用素数集 & $|S|$ & $\log_{10} N$ 下界 \\
\midrule
$S_{2012} \setminus \{398581\}$ & 14 & $3129.46$ & $S_{2012} \setminus \{19\}$ & 14 & $2144.95$ \\
$S_{2012} \setminus \{1093\}$ & 14 & $3128.17$ & $S_{2012} \setminus \{17\}$ & 14 & $2087.25$ \\
$S_{2012} \setminus \{331\}$ & 14 & $3080.59$ & $S_{2012} \setminus \{13\}$ & 14 & $1797.71$ \\
$S_{2012} \setminus \{307\}$ & 14 & $3075.50$ & $S_{\mathrm{BCR}}$（8 素数） & 8 & $1725.52$ \\
$S_{2012} \setminus \{127\}$ & 14 & $2945.55$ & $S_{2012} \setminus \{11\}$ & 14 & $1655.73$ \\
$S_{2012} \setminus \{97\}$ & 14 & $2916.71$ & $S_{2012} \setminus \{7\}$ & 14 & $1084.47$ \\
$S_{2012} \setminus \{61\}$ & 14 & $2794.52$ & $S_{2012} \setminus \{5\}$ & 14 & $784.85$ \\
$S_{2012} \setminus \{31\}$ & 14 & $2508.39$ & $S_{2012} \setminus \{3\}$ & 14 & $213.04$ \\
\bottomrule
\end{tabular}
\end{table}

表 \ref{tab:ablation} 揭示了三点关键洞察：
\begin{itemize}
    \item \textbf{8 素数追平 15 素数}：仅用 $S_{\mathrm{BCR}}$ 的 8 个素数，Model B 即可达到 $10^{1725.52}$，几乎追平 2012 年论文 15 个素数贪心法的 $10^{1735}$。
    \item \textbf{下界由前 7 个小素数主导}：剔除 $\{3, 5, 7, 11, 13, 17, 19\}$ 中任意一个素数，下界均骤降 $1000$ 个数量级以上；而剔除 $398581$、$1093$、$331$、$307$ 等较大素数几乎不影响终局下界。
    \item \textbf{对前半段的直接意义}：在冲击 $10^{2200}$ 乃至 $10^{2500}$ 无条件下界时，已无需继续向 $S$ 中追加新的难缠素数；相反，可以根据前半段因子链的实际开销，仅保留对打断小素数循环确有助益的辅助素数。
\end{itemize}

\section{路障绕过实测：Roadblock 2 案例分析}

在 Ochem 与 Rao（2012）的 Theorem 1 证明树中，总计 $22\,514\,255$ 次分支里有 $10\,406\,935$ 次（约 $46.2\%$）消耗在路障绕过上。论文图 1 给出了一个典型的双层嵌套路障案例（Roadblock 2）：
\[
M' = 3^2 \cdot 7^2 \cdot 13^1 \cdot 19^2 \cdot 127^{192} \cdot \sigma(127^{192}) > 10^{814},
\]
已知丰度上界 $a' < 1.92522$，目标规模 $B = 10^{1500}$（故剩余素因子乘积至少需超过 $10^{686}$），且已有素数集 $E = \{3, 7, 13, 19, 127\}$。原论文使用 $t_{a'}(p) = \lceil \ln(2/a') / \ln(p/(p-1)) \rceil$ 与 $p^{2t_{a'}(p)-1} > 10^{1500}/M'$，求得临界值 $b = 2625$，从而被迫对所有小于 $2625$ 且不在 $E$ 中的约 $376$ 个素数逐一分支。

将第 2 节的通用 LP 模型移植至此路障时，可同时利用两项严格结构信息：
\begin{enumerate}
    \item 唯一特殊素数已被 $13^1$ 占用，所有新引入素数只能取偶数幂次（即关闭特殊素数变量）。
    \item 设 $p$ 为整除 $N$ 但不整除 $M'$ 的最小素数，则所有满足 $q < p$ 且 $q \notin E$ 的奇素数均严格不能整除 $N$，因而全部进入有效禁用集 $S_{\mathrm{eff}}$，同时过滤掉一切 $\sigma(r^a)$ 含此类因子的分支选项。
\end{enumerate}

在设置 $P_{\max} = \max(4000, 4p)$、$a_{\mathrm{cut}} = 16$ 的条件下，对该路障执行候选素数扫描与精确二分搜索（全流程连同命题 1 对偶认证总耗时仅 $1.13$ 秒），结果见表 \ref{tab:rb2}。

\begin{table}[htbp]
\centering
\caption{Roadblock 2 候选最小新素数 $p$ 扫描与精确二分结果（剩余规模阈值 $10^{686}$）}
\label{tab:rb2}
\begin{tabular}{@{}cc|cc|cc@{}}
\toprule
最小新素数 $p$ & 剩余下界（$\log_{10}$） & 最小新素数 $p$ & 剩余下界（$\log_{10}$） & 最小新素数 $p$ & 剩余下界（$\log_{10}$） \\
\midrule
$17$ & $8.46$ & $211$ & $137.78$ & $829$ & $667.62$ \\
$53$ & $24.94$ & $307$ & $196.39$ & $853$ & $676.24$ \\
$101$ & $41.96$ & $503$ & $332.43$ & $857$ & $684.44$ \\
$131$ & $77.01$ & $739$ & $597.19$ & $\mathbf{859}$ & $\mathbf{687.10}$ \\
$167$ & $93.87$ & $809$ & $658.73$ & $991$ & $794.42$ \\
\bottomrule
\end{tabular}
\end{table}

表 \ref{tab:rb2} 表明，当最小新素数 $p \ge 859$ 时，剩余规模下界达到 $10^{687.10} > 10^{686}$，直接触发规模矛盾。因此仅需对小于 $859$ 且不在 $E$ 中的素数（共 $143$ 个）进行分支，相较原论文的 $b = 2625$（约 $376$ 个素数），单层嵌套路障的分支宽度缩减了 **$62\%$**。在多层嵌套路障中，此类宽度的压低将带来指数级的证明树剪枝收益。

\section{总结与路线 B 衔接规划}

路线 A 的阶段性工作已验证了算法层面的核心突破：在不增加任何大整数分解算力的前提下，仅通过 LP 松弛与二次因子蕴含约束，即可将 $S_{2012}$ 终局下界严格推至 $10^{3129.458561}$，并将典型嵌套路障的分支宽度砍去六成以上。

下一步向路线 B（无条件 $N > 10^{2200}+$ 前半段证明树）推进的衔接方案如下：
\begin{itemize}
    \item \textbf{构建 C++ 分解库索引与 $10^{1500}$ 基线复现}：基于 WSL2 Ubuntu 下的 \texttt{g++}、GMP 与 \texttt{zlib}，流式读取 \texttt{checkfacts.txt.gz} 建立按 $(p, q+1)$ 哈希的紧凑二进制索引，先复现 Ochem–Rao（2012）的 $10^{1500}$ 证明树以校准节点数与分支逻辑。
    \item \textbf{部署两级混合路障绕过}：在证明树运行时，先以微秒级 C++ 合法连续素数乘积界过滤浅层路障；仅当算出的素数上界仍较大（如 $b > 200$）时，再调用带二次蕴含约束的 LP 求解器（配对偶认证）压低 $b$。
    \item \textbf{路线 A 的进一步极限探索}：可将 $P_{\max}$ 扩展至 $10\,000$ 以上、引入二跳蕴含（$p \to q \to r$），并在不构成循环论证的前提下嵌入独立的最大组件下界（$p^e > 10^{62}$），进一步抬高终局冗余。
\end{itemize}

\begin{thebibliography}{9}
\bibitem{ochem2012}
P. Ochem and M. Rao, \textit{Odd perfect numbers are greater than $10^{1500}$}, Mathematics of Computation, 81 (2012), no. 279, pp. 1869--1877.

\bibitem{brent1991}
R. P. Brent, G. L. Cohen, and H. J. J. te Riele, \textit{Improved techniques for lower bounds for odd perfect numbers}, Mathematics of Computation, 57 (1991), no. 196, pp. 857--868.

\bibitem{hare2007}
K. G. Hare, \textit{New techniques for bounds on the total number of prime factors of an odd perfect number}, Mathematics of Computation, 76 (2007), no. 260, pp. 2241--2248.
\end{thebibliography}

\end{document}